\addbibresource{references.bib}

\title{Reach Men, Reach Families}
\author{Jason Reeves (July 2026)}
\date{}

\begin{document}

\ifdefined\embedversion
  \pagestyle{empty}
\fi

\maketitle
\extrafloats{100}

\begin{abstract}
The purpose of this article is to propose a ministry strategy for churches in 21st century America.  That ministry strategy involves focusing on reaching and discipling men, and prioritizing their spiritual formation with the end goal that they become true shepherds and spiritual leaders within their families.  This is not to the exclusion of ministering to women and children within the congregation -- rather, behind this proposal is a desire to work \textit{within} God's established order, to more effectively minister to the \textit{entire congregation}, but to do so in ways which comport with God's ways. 

This order was established at creation, and has been reiterated in many ways throughout the history of God's dealings with His people.  Working \textit{against} this creation order invites discord and disorder.  Working \textit{within} God's prescribed boundaries promises fruitfulness and growth -- not only to the men being enabled and equipped to properly lead their families, but also to their wives and their children to whom they have been given the task to minister.  The roles of parents in general and fathers in particular are non-transferable roles.  Well-meaning churches cannot abrogate Christian parents of their responsibilities before God, and would do well not to step in and "do the job" for the deficient parents in their midst.  Rather, the response should be equipping and discipleship of these parents so that they can better carry out their duties.  The church may assist, but may not subsume the role of parents.

We will look at the data regarding changing church demographics and influences, state the problems indicated by this data, and then propose solutions based on God's word.
\end{abstract}

\newthought{In many ways, the prosperous western church} has been on the decline since the beginning of the 20th century.  J. Gresham Machen wrote over a century ago about this decline, and his eerily prescient characterization of the problem demonstrates the relevance of his work even today.\footcite{machenchristianity}  The church has been battered by \textit{all four} waves of feminism, by various homosexual movements, transgender insanity, and, most recently, ESG\footnote{Environmental, Social, and Governance}, DEI\footnote{Diversity, Equity, and Inclusion}, and a weird principle of "safety first" in the face of global "pandemics" that many believed should override the mandate to carry out the public worship of God.  Though these most recent ideas definitely took a deeper root in corporate America than in churches, the effects of these philosophies upon the church at large cannot be denied.  Churches not scarred by embracing such ideologies (of course, \textit{contra} scripture), have been battle hardened by having to beat back these worldly philosophies.  Some churches have unfortunately \textit{over}corrected, and have been damaged as a result.

One of the ways the culture at large has been reflected in modern churches is in the overall feminization of the church.  Many observers have noted that throughout the 20th century, church architecture, design, decor, as well as worship styles and musical choices, have tended more to the feminine than the masculine.  Church programming has been engineered to suit feminine tastes.  "Soft men", many of whom are not functioning as the family shepherds God has called them to be, have been more and more preferred as elders and pastors.  The so-called Eleventh Commandment, "thou shalt be nice", manufactured though it is, has become more important than God's own Decalogue.  

Even in churches which have rejected the most flagrant violations indicated by the rebellion of the 20th century: allowing women to preach, embracing the social justice movement, affirming homosexuals in their sin, etc., the pulse and current of the culture can be felt.  Men may feel weak or wholly inadequate to shepherd their families.  Taking a cue from the culture, they may feel like being assertive and taking authority, even God-ordained authority, is somehow "toxic".  Their wives, and I say this to their shame, their very \textit{elders}, may also work to dampen their enthusiasm to simply shepherd their families as God has commanded.  This disorder is the unfortunate result of being discipled by the \textit{culture} moreso than by God's holy word.  

\ifdefined\embedversion\newpage\fi
\section{The Data}
Modern western churches overwhelmingly favor a passive version of masculinity, assuming they acknowledge masculinity at all.  As a result, men often "check out" of modern church.  Men who stay are encouraged in passivity and many don't lead their families well.  Men, who are called by God to minister to their wives and to raise their children in the faith, are replaced by women's ministries and youth ministries in a vain attempt to fill the gap, but the data shows these efforts do not achieve the intended result, because at the root what's being ignored is God's covenant order.

There\footcite{murrowmen} have\footcite{bauchamdriven} been\footcite{hamgone} many\footcite{coughlinnice} books\footcite{fosterman} written documenting the disengagement of men from modern churches -- many more than are referenced in this article.  This disengagment of men and the harmful effects upon families is nothing new.  If you lived through the \#MeToo movement in 2017, the removal of Paige Patterson from Southwestern in 2018, that horrible Gilette commercial in 2019, the COVID-19 pandemic in 2020 during which the corresponding "Summer of Love" occurred with it's "mostly peaceful" protests, and the ensuing ESG and DEI nonsense that swept the United States, then you probably witnessed some of these societal fault lines develop and fracture before your very eyes.  You may have thought your own church or denomination was immune.  Very likely, it was not.  Most churches fell into the category of total capitulation to what they perceived as the new world order, or they overcorrected -- both with negative outcomes.\footnote{Most churches that responded \textit{ad extremum} ended up dead, dying, or apostate.} Very few churches had a measured, mature response.  Of course the period of the late twenty teens to early 2020s did not mark the genesis of ideas such as "toxic masculinity" or "intersectionality" -- these ideas have been around for quite some time -- but events in the vicinity of 2020 certainly made these ideas harder to ignore, and harder to remain silent about.  Following is a survey of these references.

ALL of these books document the following:
\begin{enumerate}
  \item Churches are healthier when men are engaged.
  \item Families are healthier when men are engaged in church.
  \item Marriages are healthier when men are engaged in church.
  \item Youth and children are more likely to practice Christianity as adults when their dads are engaged in church.
  \item There is a very real gender gap -- fathers and men in general have seen a real numerical decline in churches for decades.
  \item Churches are becoming more and more feminized, creating barriers to male engagement.
  \item The church's preoccupation with creating subcultures within the church by way of programs is destructive.
\end{enumerate}

\subsection{Why Men Hate Going to Church (2005)}
David Murrow is a former Presbyterian elder who has some conclusions with which I disagree\footnote{For example, he seems to think singing is not a terribly masculine thing to do.  I think he's more than slightly bonkers here, since singing has been a core identity of God's people for millenia.  He also doesn't seem to think female pastors are part of the problem.}  However, he's right in that he's identified that the reason many men don't attend church is that the thermostat (figuratively speaking) of most churches has been set to make women and children comfortable, not to mention effete, passive men, with which most men don't want to associate. Murrow goes to great lengths to demonstrate from the data that there has been a steady decline of male participation in the western church for a century at least.  This is in spite of the reality that "religious participation leads men to become more engaged husbands and fathers.  Teens with religious fathers are more likely to say that they enjoy spending time with their dads and that they admire them.  Obviously, men need the church. \textit{But does the church need men?}"\footcite[p. 23, emphasis mine]{murrowmen}  Unfortunately the message that many men come away with is, "not really".

Murrow believes that "churches need men because men are natural risk takers"\footcite[p. 25]{murrowmen} and cites the parable of the talents to demonstrate that Christianity is a religion that must be comfortable taking risks.  Note what the position of the "wicked and lazy" servant in the parable was.  However, men's risky tendencies are a trait that critics would say is indicative of "toxic masculinity" -- it's just not safe.  But is it the case that the Lord's churches are supposed to be \textit{safe} in that way?  That the Lord's people are supposed to be \textit{safe} -- or do we not serve a God who calls us to die?

Murrow points out how liberal churches have faced steep declines in recent decades, and the decline usually starts with the men leaving.  He writes, "self-described liberal churches were 14 percent more likely to have a gender gap than conservative ones"\footcite[p. 25]{murrowmen} by which he means that there were significantly more women in such churches than men.

Churches that value stability and security over obedience, which are risk averse, which make decisions and resolve conflict the "feminine way",\footnote{On p. 33 Murrow writes, "Being a church leader is a frustrating experience, because a man cannot lead like a man.  Instead he must be careful, sentimental, and thrifty; make every decision by consensus; talk everything to death.  Decisions take months or years to make, and if someone's feelings might be hurt, we don't move forward."  Masculine men simply lose patience faced with this level of inaction.} and which insist on leadership being "nice" in all situtations, as opposed to "good" -- these types of churches do not attract the kind of men which are needed to lead those churches into the next generation.  But if your church is a church full of women, children, and effete, passive men, there may be no other options of which the congregation is aware to face such problems.  

Murrow writes that "ministries targeted at subgroups (such as women's ministry, children's ministry, youth ministry) were rare before 1800." In the Victorian era (1830s - early 1900s), clergymen realized that the newly idle \textit{nouveau riche} women whose husbands were becoming wealthy capitalists needed something to do.  The Victorian church multiplied programs, creating "ladies' teas, sewing circles, temperance unions, and women's missionary societies" in order to fill what they perceived as a ministry gap.  Children's ministry did not exist before the 1800s.  Youth ministry did not exist before the 1800s.\footcite[p. 59]{murrowmen}  Yet we cannot say that before 1800, no women, no children, and no youth were being ministered to.  \textit{How} were they being ministered to?  By the proclamation of the word, the administration of the ordinances on the Lord's day, the reading of the Bible, and family worship.  In other words, the \textit{ordinary} means of grace God has ordained for family ministry.

After World War II in the United States, the situtation became more dire.  The phenomenon of "church shopping" began, and the consumer was king.  Programs went from being nice, helpful add-ons to mandatory.  If your church didn't have a specific program being shopped for, people would just go to another church.  Murrow sadly documents, "By the 1960s these ministry programs were no longer optional: a church had to offer them, or shoppers -- excuse me; parishioners -- would flee to the competition."\footcite[p. 60]{murrowmen} Something that Murrow calls the "Christian Industrial Complex" reinforces the need for and the orientation of programs.  This is more of a market-driven than a Spirit driven phenomenon.  Publishers know what sells, and writers create what publishers want.\footnote{On p. 62, Murrow writes, "Here is the politically incorrect truth: there are measurable, verifiable differences in the way men and women percieve the world.  Whether these differences are hard-wired or socially programmed is beside the point.  They exist.  Men and women respond favorably to different things.  They're attracted to different things.  They're repelled by different things."  Knowing this, writes Murrow on the following page, "The Christian-industrial complex keeps pumping out products for women.  Christians of both genders use these products and absorb the ideas therein.  Son, everyone is looking at the Christian faith through a feminine lens.  The weight of all this female-targeted religious material is beginning to warp the faith it's supposedly describing.  The more Christian products we consume, the more we come to perceive our faith in feminine terms."}  Murrow claims that "they've bent our customs, our vocabulary, our decor, and even our theology in the feminine direction.  The raw truth is, the complex makes more money when they market to women.  So they do.  And we all pay the price.

"Toxic effeminacy," if you will, has come to dominate how we worship as well.  Murrow says "The old worship was about coming together to extol God; the new worship is about coming together to \textit{experience} God.  The target of worship has fallen half a meter -- from the head to the heart."\footcite[p. 70]{murrowmen} Effete "worship" leaders with skinny jeans and soft voices who just want to \textit{usher us into the presence} of God have become non-controversial.\footcite[p. 76]{murrowmen}

"Very well," you may say, "but my church is full of chest beating men with full beards who smoke cigars and drink bourbon."  There is absolutely nothing wrong with those things in moderation and in subjection to Christ, but they are merely \textit{outward signs}.  It is possible to have the trappings of masculinity but be blowing it at home with respect to discipling the family.  As a church we should not tolerate effeminacy, but we should also not tolerate a merely \textit{shallow} masculinity.  God calls men to be men down to their very bones, and in fact the cells of men's very bones brandish the dangerous XY set of chromosomes with which God has marked men as men.  These dangerous chromosomes are necessary to build churches, not to mention civilizations.

Ultimately, what is the answer to the question posed in Murrow's title?  Why do, in fact, many men \textit{hate} going to church? Because churches don't have anything for them to do.  Churches are neither holding them responsible nor equipping them for the ministry of male headship.  The church certainly seems to them to be a place for women, children, and effeminate or otherwise weak men.  


\subsection{Family Driven Faith (2007)}
In \textit{Family Driven Faith}, Voddie Baucham makes the case for a \textit{family integrated}\footnote{\textit{Family integrated} describes a church that does not attempt to divide the family to instruct its members separately, but rather ministers to the family as a \textit{unit}, recognizing God's ordained authority structure and operating on the family through the head, which is generally the father.  Such churches encourage families to stay together during all church activities: worship, teaching, fellowship, etc.} church. Baucham has preached\footcite{bauchamcentrality} and written and preached much on the topic of the family, most recently in \textit{Family Shepherds}\footcite{bauchamshepherds}, published in 2011.

Baucham boldly proclaims, "The family is the evangelism and discipleship arm of the family-integrated church.  This is the aspect of the family-integrated paradigm that critics find most disturbing."\footcite[p. 195]{bauchamdriven} Why do they find it disturbing?  Baucham highlights that many churches and pastors simply don't trust parents to evangelize and disciple their kids.  Many parents don't even trust \textit{themselves} to evangelize and disciple their kids.  He continues on the next page, "There seems to be a general consensus that it is unreasonable to expect families to function the way the Bible clearly says they should."\footcite[p. 196]{bauchamdriven}  Baucham documents various "gimmicks" that churches have tried to use to attract men.  Baucham concludes, "Ironically, these approaches all leave out the preeminent biblical mandate that God has given to men -- the mandate to love their wives as Christ loved the church and to bring their children up in the nurture and admonition of the Lord."\footnote{Ephesians 5:25 - 6:4} He claims that because churches fail to hold men accountable, because churches fail to expect much of their men, and are reticent to do the hard work to \textit{equip} even the willing men among them to do the job God has called them to do.  

Baucham spends much time in this book and in many of his videos decrying his observations about a "youth ministry subculture" that creates a "church within a church" and does not put young people on the road to becoming mature, adult followers of Christ.  He contends repeatedly that it is the job of parents to raise children, primarily the father.  The statistics point out that between 70 and 88\% of children and youth who attend age segregated services in churches are not serving the Lord as adults.\footcite[p. 10]{bauchamdriven}  If parents, primarily fathers, don't show up and demonstrate to their children that 1) Churches are \textit{not} just places for women, children, and weak men, and that 2) their fathers (preferably fathers AND mothers) are heavily engaged in the faith, both mentally and physically, and with regard to the very structure of their lives, then these children show a strong preference to walk away from the faith.\footnote{They went out from us, but they were not really of us; for if they had been of us, they would have remained with us; but they went out, in order that it might be shown that they all are not of us. \scriptref{I John 2:19-20 \biblever{nasb 77}}} I believe also worth considering -- if children languish spiritually in youth ministries which operate outside of the family model, how do women fare with regard to women's ministries which also are outside of the family model?  Now there are always outliers, but ask yourself, do women's ministries \textit{generally} result in more holy living?  Better marriages?  More Christlike moms?  "Well, I can't depend on my husband to disciple me!"  Perhaps not.  But that's the \textit{biblical}\footnote{Ephesians 5:22-33 is instructive here, as well as I Corinthians 14:34-35.  In my opinion, the mandate for husbands to disciple wives is hard to miss.}  model.  This question of women's ministries is not one that Baucham addresses in his book, but seems to arise naturally out of the conversation -- if working outside of God's ordained authority structure is bad for children, it follows that there might be deleterious effects for everyone.  

Baucham points out that Christian parents are opting not to have sufficiently large families, and when they do, they are doing a poor job of passing the faith to their children.\footcite[p. 174]{bauchamdriven}  Baucham notes, "With a birth rate hovering around two children per family, a biblical worldview rate below 10 percent, and about 75 percent of our teens leaving the church by the end of their freshman year in college, \textit{it currently takes two Christian families in one generation to get a single Christian into the next generation.}"\footcite[p. 174, emphasis mine.]{bauchamdriven}  These statistics are unsustainable, and the implications are dire.  Are we willing to give up our comfort in order to reconfigure church ministries so that we are doing things God's way, a way He promises to bless?\footnote{Deuteronomy 28:1-6}

\subsection{Already Gone (2009)}
Ken Ham in \textit{Already Gone} also points out the challenging statistics facing Christians today.  Ham quotes George Barna, "A majority of twenty-somethings -- 61\% of today's young adults -- had been churched at one point during their teen years but they are not spiritually disengaged (i.e., not actively attending church, reading the Bible, or praying)."\footcite[p. 21]{hamgone}  Since this book is not exactly a \textit{new} work, I'm assuming the situation is much worse now than it was then.  

Ham makes the shocking claim that many of the children we see active in churches today are basically "dead men walking".  Ham notes that "six out of ten 20-somethings who were involved in a church during their teen years are already gone.  Despite strong levels of spiritual activity during the teen years, most 20-somethings disengage from active participation in the Christian faith during their young adult years -- and often beyond that."\footcite[pp. 22-23]{hamgone}   Interestingly, and tragically, many of these young adults had \textit{already made} the decision to deny the faith while in the midst of participating in these church activities.\footcite[p. 32]{hamgone}

Ham points out two of the same concerns that Baucham identified - 1) churches encourage, and parents happily engage in outsourcing their parental duties\footcite[pp. 50-51]{hamgone}, and 2) churches neglect to equip parents to take on the task God has given, non-transferably, to them.\footcite[p. 154]{hamgone}

\subsection{No More Christian Nice Guy (2016)}

\subsection{It's Good to Be A Man (2021)}

\ifdefined\embedversion\newpage\fi
\section{The Problem}
Churches err when they try to circumvent God's covenant order as well as violate the concept of sphere sovereignty.  The Bible teaches male headship. Although men and women are of equal worth before God, they have different callings based upon immutable biological and spiritual characteristics.\footnote{This biblical way of understanding gender roles and differences is called \textit{complementarianism}, meaning to indicate that men and women were created to \textit{complement} each other, each making up for the inherent deficiencies in the other.  Throughout this discussion, I'll be using \textit{gender} and \textit{sex} as synonyms recognizing immutable biological and spiritual characteristics, not social constructs.  The notion that one's God-given gender is merely a "social construct" is anathema.}  The role God has chosen for men is that of the head, or leader within families, churches, and indeed civilization itself.  We have to work \textit{within} the frameworks God has provided for us, or else we risk distorting God's truth, not to mention doing harm to ourselves and others.

\ifdefined\embedversion\newpage\fi
\section{The Solution}
Churches should focus on and prioritize ministry to men in order to build strong families and strong churches.  Reaching men reaches everyone.  By working \textit{along with} God's design, we will be able to do His work more effectively.  It's already hard enough to do God's work in this world, but we make it harder when we "kick against the pricks".  
\end{document}
